% !TeX program = xelatex
% ==========================================================================
% Пропоузал семестрового проєкту — курс «Штучний інтелект», УКУ.
% Це ЄДИНИЙ файл, який вам потрібно редагувати. Уся верстка й брендинг
% живуть у ucuproposal.cls — не змінюйте шрифти, поля, кеглі й кольори.
%
% Компіляція: XeLaTeX + Biber (Overleaf робить це автоматично).
% Локально:   latexmk (див. README.md).
% ==========================================================================
\documentclass[ukrainian]{ucuproposal}
% Опції класу (можна комбінувати через кому):
%   ukrainian | english  — мова документа й логотипу (за замовчуванням ukrainian)
%   slogan                — додає гасло УКУ в колонтитул титулки (за замовчуванням вимкнено)

% --------------------------------------------------------------------------
% Метадані титульної сторінки
% --------------------------------------------------------------------------
\projecttitle{Децентралізований рій роботів Codey Rocky}


\teammember{Дзюба Оксана}
\teammember{Сова Юлія}
% \teammember{Прізвище Ім'я} % третій учасник (опційно)
% \teammember[https://linkedin.com/in/...]{Прізвище Ім'я} % опційне клікабельне посилання на профіль (LinkedIn тощо)

\mentor{} % якщо ментора немає — видаліть текст і залиште \mentor{}, рядок не друкується

\track{research} % research | engineering — не друкується на титулці, лише для внутрішньої категоризації

\repo{https://github.com/JuliaSovaO/CATaclysm}

\proposaldate{2026-09-21}

\begin{document}
\maketitlepage

\section{Проблема та мотивація}
Рій — це група агентів, що досягає складної колективної поведінки, приймаючи рішення локально й незалежно, без централізованого керування. Це один із ключових напрямків сучасного ШІ (multi-agent systems, decentralized control) із прямим застосуванням у пошуково-рятувальних операціях, моніторингу територій та безпілотних системах. Основна мета проєкту — навчити рій фізичних роботів Codey Rocky координованій поведінці за допомогою мультиагентного навчання з підкріпленням (MARL) у симуляції, з подальшим перенесенням навченої політики на реальні роботи (sim-to-real). Кожен робот діє автономно, приймаючи рішення виключно на основі власних локальних сенсорних даних у реальному часі (наявність цілі, найближчої перешкоди, сусіда в полі зору), без прямого обміну командами між роботами чи центрального координатора. Це дозволяє дослідити, чи виникає узгоджена групова поведінка (наприклад, спільне оточення цілі) як результат навченої, а не заданої вручну, локальної політики кожного агента.

\section{Огляд джерел та наявних рішень}
Базовою моделлю децентралізованої групової поведінки є класична робота Рейнольдса про flocking (правила separation/alignment/cohesion)~\cite{reynolds1987flocks}; систематичний огляд підходів до дизайну поведінки роїв роботів наведено у~\cite{brambilla2013swarm}. Ми відходимо від підходу з вручну заданими правилами на користь навчання поведінки: \cite{lowe2017maddpg} MADDPG — алгоритм мультиагентного actor-critic навчання з підкріпленням, придатний саме для змішаних кооперативно-конкурентних сценаріїв (переслідування, оточення цілі), який ми плануємо використати як основний метод. Як середовище для тренування використовуємо \textsc{PettingZoo}~\cite{terry2021pettingzoo} — стандартизовану бібліотеку мультиагентних середовищ (сценарій \texttt{simple\_tag}, predator--prey), що дозволяє задавати локальні, обмежені спостереження для кожного агента, необхідні для коректного sim-to-real переносу.

\section{Дані}
На відміну від задач із класифікацією чи детекцією зображень, дані в цьому проєкті — це не завантажений зовнішній датасет, а \textbf{симульований досвід взаємодії агентів}, що генерується під час MARL-тренування в середовищі PettingZoo (\texttt{simple\_tag}): мільйони кроків (стан $\rightarrow$ дія $\rightarrow$ нагорода), зібраних агентами в процесі проб і помилок. Ключова умова коректності підходу — спостереження кожного агента обмежені \textbf{локальними сенсорними даними в реальному часі} (відстань і напрямок до найближчої перешкоди, сусіда чи цілі в межах радіусу датчика), без доступу до глобальної карти середовища. Розташування перешкод, кількість перешкод і початкові позиції агентів \textbf{рандомізуються в кожному епізоді тренування}, а частина епізодів включає рухомі перешкоди — це не дає політиці завчити конкретну карту й змушує її навчитися загальному реактивному правилу, що коректно переноситься на фізичні роботи (Codey Rocky також не має доступу до карти кімнати, лише до миттєвих показань власних дистанційних датчиків). Додатково для оцінки якості sim-to-real переносу збираємо власні \textbf{сенсорні логи} з реальних роботів під час фізичних експериментів (дистанційні показники, прийняті рішення) — зберігаються у \texttt{data/} репозиторію. Персональних даних не передбачається.

\section{Підхід і методи}
\textbf{Основна ціль.} Тренуємо мультиагентну політику координації рою алгоритмом MADDPG (або, як спрощена альтернатива, незалежний PPO на кожного агента) у середовищі PettingZoo (\texttt{simple\_tag}). Спостереження кожного агента — виключно локальні: відстань/напрямок до найближчої перешкоди, найближчого сусіда-робота та цілі в межах обмеженого радіусу датчика (без глобальної карти). Навчена політика переноситься на фізичні роботи Codey Rocky (sim-to-real): або спрощується/дистилюється у форму, яка виконується безпосередньо на мікроконтролері, або inference виконується на зовнішньому пристрої, що керує роботом по Bluetooth/Wi-Fi.

\textbf{Додаткові цілі (за наявності часу й ресурсів), у порядку пріоритету:}
\begin{enumerate}
  \item \emph{Обхід динамічних перешкод} — розширення тренувального середовища: перешкоди рандомізуються за кількістю й позицією в кожному епізоді, частина епізодів включає рухомі перешкоди, щоб політика навчилась реагувати на сенсорний сигнал у моменті, а не на завчену розкладку.
  \item \emph{Переслідування ``розумної'' жертви} — ціль (\texttt{prey}) також керується навченим агентом, що намагається втекти (adversarial predator--prey), замість статичної чи випадково рухомої цілі.
  \item \emph{Диференційована поведінка залежно від ситуації} — аналіз навченої політики на предмет виникнення різних стратегій залежно від того, чи ціль видно, чи поруч інший робот рою, чи кілька роботів одночасно бачать ціль.
\end{enumerate}

\section{Оцінювання результатів}
\textbf{Кількісні метрики (в симуляції):} середня нагорода за епізод, частка успішних сценаріїв оточення/захоплення цілі, середня кількість зіткнень з перешкодами за епізод, швидкість збіжності навчання (кількість епізодів до стабільного результату). \textbf{Кількісні метрики (на реальних роботах):} частка успішних сценаріїв оточення цілі за фіксований час, кількість зіткнень з перешкодами за прогін, розбіжність поведінки між симуляцією та реальністю (sim-to-real gap). \textbf{Baseline.} (1) один агент без координації з рештою рою; (2) рій із вручну заданими реактивними правилами (Boids-стиль) замість навченої політики — порівнюємо навчений підхід із обома базовими варіантами. \textbf{Якісна оцінка:} відеозаписи прогонів у симуляції та на реальних роботах, що демонструють наявність або відсутність узгодженої координації.

\section{План і ризики}
До попереднього захисту: реалізація й налаштування середовища PettingZoo з локальними спостереженнями та рандомізованими перешкодами; перше успішне тренування MARL-політики базової координації рою в симуляції; порівняння з ручним (Boids) baseline. До фінального захисту: перенесення навченої політики на фізичні роботи Codey Rocky (sim-to-real), оцінка розбіжності з симуляцією, за наявності часу — динамічні перешкоди та/або переслідування рухомої цілі, повний прогін рою зі збором метрик і відео, постер та фінальна документація.

\textbf{Ключові ризики та план Б:}
\begin{itemize}
  \item MARL-тренування нестабільне або не збігається у відведений час $\rightarrow$ спрощення до незалежного PPO на кожного агента замість MADDPG, або зменшення розмірності спостережень.
  \item Sim-to-real gap виявиться завеликим (навчена політика не переноситься на реальні сенсори Codey Rocky) $\rightarrow$ калібрування симульованих сенсорів під реальні характеристики датчиків, або дистиляція навченої політики у прості реактивні правила.
  \item Обмежена обчислювальна потужність Codey Rocky для inference на борту $\rightarrow$ виконання політики на зовнішньому пристрої з передачею команд роботу по Bluetooth/Wi-Fi.
  \item Координація рою не виникає в очікуваній формі $\rightarrow$ чесно задокументувати частковий/негативний результат і проаналізувати причини (це прямо дозволено критеріями курсу).
\end{itemize}

\begin{keyfacts}
\textbf{Ключові цифри:} 2--4 роботи Codey Rocky в рої \textperiodcentered\ 1 MARL-політика (MADDPG) \textperiodcentered\ 4 тижні до midterm
\end{keyfacts}

\section{Contribution statement}
% Хто за що відповідає. Відсотки не обов'язкові — якщо хочете їх вказати,
% додайте в квадратних дужках: \contribution[50\%]{Прізвище Ім'я}{...}
\begin{contributions}
  \contribution[50\%]{Сова Юлія}{скінченний автомат станів, координація рою, документація}
  \contribution[50\%]{Дзюба Оксана}{модуль переслідування цілі,  обхід перешкод, збір метрик}
\end{contributions}

% Примітка: декларація використання AI-інструментів (AI usage) до цього
% пропоузалу не входить — вона живе в README.md вашого репозиторію.

\ucuprintbibliography

\end{document}
